% ============================================================
%  EXAMEN PARCIALITO 5  —  Métodos Numéricos
%  Compilar:  make   (o  latexmk -pdf 5toparcialito.tex)
% ============================================================
\documentclass[12pt,letterpaper]{article}

% ── Codificación e idioma ──────────────────────────────────
\usepackage[utf8]{inputenc}
\usepackage[T1]{fontenc}
\usepackage[spanish]{babel}

% ── Matemáticas ───────────────────────────────────────────
\usepackage{amsmath,amssymb,mathtools}

% ── Tablas ────────────────────────────────────────────────
\usepackage{booktabs,array,multirow}

% ── Listas ────────────────────────────────────────────────
\usepackage{enumitem}

% ── Figuras ───────────────────────────────────────────────
\usepackage{graphicx}
\graphicspath{{./capturas/}}   % imágenes en subcarpeta capturas/

% ── Geometría ─────────────────────────────────────────────
\usepackage[top=2.2cm,bottom=2.2cm,left=2.5cm,right=2.5cm]{geometry}

% ── Colores y cajas ───────────────────────────────────────
\usepackage{xcolor,mdframed}

% ── Encabezado / pie de página ────────────────────────────
\usepackage{fancyhdr}
\pagestyle{fancy}
\fancyhf{}
\rhead{\small Métodos Numéricos}
\lhead{\small Francisco Javier Ponce Sáenz — a325000}
\cfoot{\thepage}
\renewcommand{\headrulewidth}{0.4pt}

% ── Contadores de problemas ───────────────────────────────
\newcounter{problema}
\newenvironment{problema}[1][]{%
  \refstepcounter{problema}%
  \begin{mdframed}[
      linewidth=1pt,
      linecolor=black!40,
      backgroundcolor=gray!6,
      innertopmargin=6pt,
      innerbottommargin=6pt
  ]
  \noindent\textbf{Problema \theproblema}%
  \ifx&#1&\else\ \textbf{— #1}\fi\par\smallskip
}{%
  \end{mdframed}\vspace{6pt}
}

\newenvironment{solucion}{%
  \vspace{4pt}\noindent\textbf{Solución:}\par\smallskip
}{%
  \par
}

% ── Comandos útiles ───────────────────────────────────────
\newcommand{\respuesta}[1]{\par\medskip\noindent\fbox{\textbf{Resultado:}~#1}}
\newcommand{\tol}[1]{\textrm{Tol} = #1}

% ─────────────────────────────────────────────────────────
\begin{document}
% ─────────────────────────────────────────────────────────

% ══════════════════  PORTADA  ════════════════════════════
\begin{center}
  {\large\bfseries Universidad Autónoma de Chihuahua}\\[2pt]
  {\large Facultad de Ingeniería}\\[6pt]
  \rule{\linewidth}{0.8pt}\\[4pt]
  {\LARGE\bfseries Examen Parcialito 5}\\[4pt]
  \rule{\linewidth}{0.8pt}\\[6pt]
  {\large Métodos Numéricos}\\[4pt]
  {\large Francisco Javier Ponce Sáenz \quad|\quad a325000}\\[2pt]
  {\large\today}
\end{center}

\vspace{8pt}

% ══════════════════  PROBLEMA 1  ════════════════════════
\begin{problema}[Ajuste de Curvas por Mínimos Cuadrados]

De la siguiente serie de puntos:

\begin{center}
\begin{tabular}{c|rrrrrrr}
  \toprule
  $x$ & $-4$ & $-3$ & $-2$ & $-1$ & $0$ & $1$ & $2$ \\
  \midrule
  $y$ & $-3$ & $11$ & $13$ & $9$ & $5$ & $7$ & $21$ \\
  \bottomrule
\end{tabular}
\end{center}

Obtener:
\begin{enumerate}[label=\alph*)]
  \item Ajuste de curvas de 1er orden.
  \item Ajuste de curvas de 2do orden.
  \item Ajuste de curvas de 3er orden.
  \item ¿Cuál de las tres funciones es la más exacta y por qué?
  \item Completar la tabla para $x \in \{-5,\,3,\,4,\,5\}$.
\end{enumerate}

\noindent\textit{Notas: (1) El SELS formado se resuelve por eliminación de Gauss.
(2) Se incluye el procedimiento de elaboración del SELS y su resolución.}

% ───── SOLUCIÓN ──────────────────────────────────────────
\begin{solucion}

% ── Sumas ─────────────────────────────────────────────
\subsection*{Sumas necesarias ($n = 7$)}

\[
  \sum x = -7, \quad
  \sum y = 63, \quad
  \sum x^2 = 35, \quad
  \sum x^3 = -91, \quad
  \sum x^4 = 371
\]
\[
  \sum x^5 = -1267, \quad
  \sum x^6 = 4955, \quad
  \sum xy = -7, \quad
  \sum x^2 y = 203, \quad
  \sum x^3 y = -43
\]

\begin{center}
  \includegraphics[width=0.82\linewidth]{08_sistema_normal.png}\\
  {\small\itshape Sumas necesarias y sistema de ecuaciones normales (grado 3).}
\end{center}

% ── Inciso a ──────────────────────────────────────────
\subsection*{a) Ajuste de 1er orden — $Y_1 = K_0 + K_1 x$}

Sistema normal (2×2):
\[
  \begin{pmatrix} 7 & -7 \\ -7 & 35 \end{pmatrix}
  \begin{pmatrix} K_0 \\ K_1 \end{pmatrix}
  =
  \begin{pmatrix} 63 \\ -7 \end{pmatrix}
\]

Eliminación de Gauss: $R_2 \leftarrow R_2 + R_1$ y $R_2 \leftarrow R_2/28$:
\[
  K_1 = 2, \qquad K_0 = 11
\]

\[
  \boxed{Y_1 = 11 + 2x} \qquad \text{RSS}_1 = 216
\]

% ── Inciso b ──────────────────────────────────────────
\subsection*{b) Ajuste de 2do orden — $Y_2 = K_0 + K_1 x + K_2 x^2$}

Al resolver el sistema 3$\times$3 por Gauss se llega a $84\,K_2 = 0$,
por lo que $K_2 = 0$ y los demás coeficientes quedan igual que en $Y_1$:

\[
  K_0 = 11, \quad K_1 = 2, \quad K_2 = 0
\]

\begin{mdframed}[linecolor=orange!80!black, linewidth=1.2pt, backgroundcolor=orange!8]
Como $K_2 = 0$, $Y_2$ es igual a $Y_1$. El ajuste de 2do orden no aplica para estos datos.
\[
  Y_2 = 11 + 2x + 0\cdot x^2 = Y_1 \qquad \text{RSS}_2 = 216
\]
\end{mdframed}

% ── Inciso c ──────────────────────────────────────────
\subsection*{c) Ajuste de 3er orden — $Y_3 = K_0 + K_1 x + K_2 x^2 + K_3 x^3$}

Sistema normal (SSELS 4×4):
\[
  \begin{pmatrix}
    7   & -7   & 35   & -91  \\
   -7   & 35   & -91  & 371  \\
   35   & -91  & 371  & -1267 \\
   -91  & 371  & -1267 & 4955
  \end{pmatrix}
  \begin{pmatrix} K_0 \\ K_1 \\ K_2 \\ K_3 \end{pmatrix}
  =
  \begin{pmatrix} 63 \\ -7 \\ 203 \\ -43 \end{pmatrix}
\]

Resuelto por Gauss-Jordan $\Rightarrow$
$K_0 = 5,\; K_1 = -2,\; K_2 = 3,\; K_3 = 1$.

\[
  \boxed{Y_3 = 5 - 2x + 3x^2 + x^3} \qquad \text{RSS}_3 = 0
\]

% ── Inciso d ──────────────────────────────────────────
\subsection*{d) Función más exacta}

Dado que $K_2 = 0$, el ajuste de 2do orden ($Y_2$) es idéntico al de 1er orden
($Y_1$) y \textbf{se descarta como función independiente}.
La comparación real es entre $Y_1$ y $Y_3$:

\begin{center}
\begin{tabular}{lcc}
  \toprule
  Función & Válida & RSS \\
  \midrule
  $Y_1 = 11 + 2x$ & Sí & 216 \\
  $Y_2 = 11 + 2x + 0\,x^2$ & \textbf{No} ($K_2=0$, $Y_2\equiv Y_1$) & 216 \\
  $Y_3 = 5 - 2x + 3x^2 + x^3$ & Sí & \textbf{0} \\
  \bottomrule
\end{tabular}
\end{center}

Se elige $\mathbf{Y_3}$ porque $\text{RSS}_3 = 0$: el polinomio de grado 3
pasa \emph{exactamente} por los siete puntos de la tabla.

% ── Inciso e ──────────────────────────────────────────
\subsection*{e) Predicciones con $Y_3$}

\begin{center}
\begin{tabular}{r|r}
  \toprule
  $x$ & $Y_3(x) = 5 - 2x + 3x^2 + x^3$ \\
  \midrule
  $-5$ & $\mathbf{-35}$ \\
  $3$  & $\mathbf{53}$  \\
  $4$  & $\mathbf{109}$ \\
  $5$  & $\mathbf{195}$ \\
  \bottomrule
\end{tabular}
\end{center}

\respuesta{$Y_3 = 5 - 2x + 3x^2 + x^3$\quad (RSS = 0; ajuste exacto de 3er orden)}

\end{solucion}
\end{problema}

% ══════════════════  EVIDENCIA COMPUTACIONAL  ════════════════════
\newpage
\section*{Evidencia computacional — Hoja de cálculo}

Las siguientes capturas corresponden a los cuadros generados por el
programa \texttt{ajuste\_curvas\_app.py} al ejecutar la hoja de cálculo
con los datos del problema.

% ── Figura 1 ─────────────────────────────────────────────
\begin{figure}[h!]
  \centering
  \includegraphics[width=0.45\linewidth]{01_datos.png}
  \caption{Tab \emph{Datos de entrada} — tabla de los 7 puntos.}
\end{figure}

% ── Figura 2 ─────────────────────────────────────────────
\begin{figure}[h!]
  \centering
  \includegraphics[width=\linewidth]{02_procedimiento.png}
  \caption{Tab \emph{Tabla del procedimiento} — potencias de $x$, productos y sumas.}
\end{figure}

\newpage

% ── Figura 3 ─────────────────────────────────────────────
\begin{figure}[h!]
  \centering
  \includegraphics[width=\linewidth]{03_desviaciones.png}
  \caption{Tab \emph{Tabla de desviaciones} — La columna $Y_2$ (naranja) es idéntica a $Y_1$ porque $K_2=0$.}
\end{figure}

% ── Figura 4 ─────────────────────────────────────────────
\begin{figure}[h!]
  \centering
  \includegraphics[width=0.88\linewidth]{04_grafica.png}
  \caption{Tab \emph{Gráfica} — curvas de ajuste superpuestas con los datos originales.}
\end{figure}

\newpage

% ── Figura 5 ─────────────────────────────────────────────
\begin{figure}[h!]
  \centering
  \includegraphics[width=0.50\linewidth]{05_evaluacion.png}
  \caption{Tab \emph{Evaluación exacta} — $Y_3$ evaluada en $x \in \{-5, 3, 4, 5\}$.}
\end{figure}

% ── Figura 6 ─────────────────────────────────────────────
\begin{figure}[h!]
  \centering
  \includegraphics[width=\linewidth]{07_ssels.png}
  \caption{Tab \emph{SSELS} — eliminación de Gauss-Jordan del sistema 4×4 paso a paso.}
\end{figure}

% ══════════════════  FIN  ═══════════════════════════════
\end{document}
